\section{Introdução}

A programação assistida por modelos de linguagem mudou de natureza em poucos anos. O que começou como sugestão de trechos de código a partir de prompts isolados evoluiu para sistemas que decompõem o ciclo de desenvolvimento entre agentes especializados, mantêm artefatos persistentes, executam ferramentas e validam o próprio trabalho. ChatDev e MetaGPT mostraram que organizar agentes por papéis de engenharia, e não apenas melhorar o prompt, muda o que o sistema consegue entregar \cite{ChenQian2023,SiruiHong2023}. A literatura mais recente já fala em uma Engenharia de Software nativa de IA e trata sistemas multiagente como linha de pesquisa própria \cite{hou2024llmse,AhmedEHassan2024,JundaHe2025}. A unidade de trabalho, antes o prompt, tornou-se o processo.

Esse deslocamento traz um objeto novo: o framework operacional que organiza o desenvolvimento sobre um agente. Caracterizá-lo exige examinar a função de cada mecanismo, como a especificação dirige a geração, como o contexto do repositório é recuperado, como papéis são distribuídos, como a execução age sobre o ambiente e como a validação ancora o resultado, em vez de métricas superficiais como número de comandos ou popularidade. As revisões existentes, porém, abordam o problema por outro ângulo. Parte da literatura classifica arquiteturas de agentes de forma abstrata, por componentes cognitivos, e parte mapeia técnicas de modelos de linguagem por tarefas isoladas do ciclo de vida. Falta uma lente de Engenharia de Software voltada para os frameworks que organizam o ciclo de desenvolvimento em nível de processo e de repositório, e falta confrontar a estrutura desses frameworks com a evidência disponível, em vez de descrevê-los um a um.

Para preencher essa lacuna, conduzimos uma revisão sistemática de literatura sob o protocolo de Kitchenham, restrita a artigos acadêmicos e preprints arquivados com registro verificável \cite{kitchenham2009slr}. A partir de um corpus final de 41 estudos publicados entre 2022 e 2026, propomos uma taxonomia de seis dimensões para classificar esses frameworks e a confrontamos com a evidência do corpus, por meio de uma síntese temática em modo híbrido. As contribuições deste artigo são três: uma taxonomia de seis dimensões (especificação, contexto, papéis, execução, validação e portabilidade) testada contra a evidência, e não assumida a priori; uma síntese temática de seis temas indutivos que cobre 100\% do corpus e expõe três eixos de tensão e dois eixos emergentes fora da taxonomia; e uma agenda de pesquisa orientada a processo que decorre diretamente das lacunas observadas.

A revisão é guiada por cinco perguntas de pesquisa. RQ1: quais dimensões arquiteturais distinguem esses frameworks? RQ2: como eles organizam contexto, papéis, artefatos, execução e validação? RQ3: que tipos de evidência sustentam as alegações? RQ4: quais riscos e limitações são reportados ou inferíveis? RQ5: que lacunas de pesquisa permanecem para comparar, avaliar e governar esses frameworks em projetos reais? O restante do artigo está organizado da seguinte forma: o Método de Revisão descreve o protocolo e o fluxo de seleção; os Trabalhos Relacionados posicionam a revisão frente à literatura existente; a Taxonomia Proposta apresenta as seis dimensões; os Resultados e Síntese consolidam os seis temas e confrontam a taxonomia com a evidência; a Discussão expõe os achados transversais e os eixos de tensão; e a Agenda de Pesquisa e a Conclusão delineiam os próximos passos.
